% step(z), the threshold the linear binary classifier applies to z.
%
% Used on lecture 4 badges 6, 9 and 13, which all share one raster. Black and
% white on purpose: the plot is one curve, so colour would only decorate.
%
% The label is \operatorname{step}, not Step or $Step$. Writing $Step(z)$ sets
% S, t, e and p as four italic variables in a product, which is what the raster
% this replaces used to show. \operatorname also puts the correct thin space
% before the parenthesis, and matches the deck's own math on badge 9,
% g = \operatorname{step}(z) = \operatorname{step}(\theta^\top x + \theta_0).
%
% Build:  pdflatex step.tex  &&  ../pdf2png.sh step

\documentclass[tikz]{standalone}
\usepackage{pgfplots}
\usepackage{amsmath}
\pgfplotsset{compat=1.18}

\begin{document}
\begin{tikzpicture}
  \begin{axis}[
    axis lines=middle,
    x axis line style={->,>=stealth, shorten >=-3pt},
    y axis line style={->,>=stealth, shorten >=-3pt},
    xmin=-4.4, xmax=4.4,
    % ymin at 0 so the y axis has no tail below the origin. The lower half of
    % the step then lies along the x axis, which is what it should look like.
    ymin=0, ymax=1.42,
    xtick={-4,-2,2,4},
    ytick={1},
    extra x ticks={0}, extra x tick style={xticklabel style={anchor=north east}},
    tick style={black, thin},
    axis line style={black},
    width=10cm, height=4.75cm,
    scale only axis,
    clip=false,
  ]
    % the guide to 1, drawn before the curve so the curve sits on top
    \addplot[gray, dashed, thin, domain=-4:0, samples=2] {1};

    % step(z): 0 below the threshold, 1 above it, with the jump at z = 0.
    % Three separate plots rather than one sampled curve, so the vertical is
    % exactly vertical instead of a steep sampled ramp.
    \addplot[black, line width=2pt, domain=-4:0,  samples=2] {0};
    \addplot[black, line width=2pt] coordinates {(0,0) (0,1)};
    \addplot[black, line width=2pt, domain=0:4,   samples=2] {1};

    % Axis names as nodes rather than xlabel/ylabel: both sit beside their
    % arrow tip, which pgfplots' label anchoring will not do on a middle axis.
    \node[anchor=south west] at (axis cs:0.12,1.20) {\(\operatorname{step}(z)\)};
    \node[anchor=south east] at (axis cs:4.38,0.05) {\(z\)};
  \end{axis}
  \pgfresetboundingbox
  \useasboundingbox ([xshift=-9mm, yshift=-7mm]current axis.south west)
    rectangle ([xshift=6mm, yshift=6mm]current axis.north east);
\end{tikzpicture}
\end{document}
